\documentclass[10pt,a4paper]{amsart}
\usepackage{amsmath,amssymb,amsthm,graphicx}
% Retained minimal preamble for the Stacks Project chapter more-algebra.tex (More on Algebra).
% The chapter file pinned-original-source.tex carries only \input{preamble} on its first line, so no
% preamble text lives inside the chapter itself.  The Stacks Project uses one single shared preamble for
% every chapter; the snapshot stages it identically as sources/stacks-chow/preamble.tex and
% sources/stacks-etale-cohomology/preamble.tex.  The block below is copied mechanically, line for line,
% from sources/stacks-chow/preamble.tex lines 14-18 (the Xy-pic setup) and 154-217 (the native theorem
% declarations together with the mathematical macros).  Nothing was added, removed or redefined; the
% wrapper class and the amsmath/amssymb/amsthm/graphicx base packages are supplied by the assembler.
% Chapter: more-algebra.tex -- More on Algebra.
% For commutative diagrams we use Xy-pic
\usepackage[all]{xy}

% We use 2cell for 2-commutative diagrams.
\xyoption{2cell}

% Theorem environments.
%
\theoremstyle{plain}
\newtheorem{theorem}[subsection]{Theorem}
\newtheorem{proposition}[subsection]{Proposition}
\newtheorem{lemma}[subsection]{Lemma}

\theoremstyle{definition}
\newtheorem{definition}[subsection]{Definition}
\newtheorem{example}[subsection]{Example}
\newtheorem{exercise}[subsection]{Exercise}
\newtheorem{situation}[subsection]{Situation}

\theoremstyle{remark}
\newtheorem{remark}[subsection]{Remark}
\newtheorem{remarks}[subsection]{Remarks}

\numberwithin{equation}{subsection}

% Macros
%
\def\lim{\mathop{\mathrm{lim}}\nolimits}
\def\colim{\mathop{\mathrm{colim}}\nolimits}
\def\Spec{\mathop{\mathrm{Spec}}}
\def\Hom{\mathop{\mathrm{Hom}}\nolimits}
\def\Ext{\mathop{\mathrm{Ext}}\nolimits}
\def\SheafHom{\mathop{\mathcal{H}\!\mathit{om}}\nolimits}
\def\SheafExt{\mathop{\mathcal{E}\!\mathit{xt}}\nolimits}
\def\Sch{\mathit{Sch}}
\def\Mor{\mathop{\mathrm{Mor}}\nolimits}
\def\Ob{\mathop{\mathrm{Ob}}\nolimits}
\def\Sh{\mathop{\mathit{Sh}}\nolimits}
\def\NL{\mathop{N\!L}\nolimits}
\def\CH{\mathop{\mathrm{CH}}\nolimits}
\def\proetale{{pro\text{-}\acute{e}tale}}
\def\etale{{\acute{e}tale}}
\def\QCoh{\mathit{QCoh}}
\def\Ker{\mathop{\mathrm{Ker}}}
\def\Im{\mathop{\mathrm{Im}}}
\def\Coker{\mathop{\mathrm{Coker}}}
\def\Coim{\mathop{\mathrm{Coim}}}

% Boxtimes
%
\DeclareMathSymbol{\boxtimes}{\mathbin}{AMSa}{"02}

%
% Macros for moduli stacks/spaces
%
\def\QCohstack{\mathcal{QC}\!\mathit{oh}}
\def\Cohstack{\mathcal{C}\!\mathit{oh}}
\def\Spacesstack{\mathcal{S}\!\mathit{paces}}
\def\Quotfunctor{\mathrm{Quot}}
\def\Hilbfunctor{\mathrm{Hilb}}
\def\Curvesstack{\mathcal{C}\!\mathit{urves}}
\def\Polarizedstack{\mathcal{P}\!\mathit{olarized}}
\def\Complexesstack{\mathcal{C}\!\mathit{omplexes}}
% \Pic is the operator that assigns to X its picard group, usage \Pic(X)
% \Picardstack_{X/B} denotes the Picard stack of X over B
% \Picardfunctor_{X/B} denotes the Picard functor of X over B
\def\Pic{\mathop{\mathrm{Pic}}\nolimits}
\def\Picardstack{\mathcal{P}\!\mathit{ic}}
\def\Picardfunctor{\mathrm{Pic}}
\def\Deformationcategory{\mathcal{D}\!\mathit{ef}}

% --- appended 2026-10-12 (Chief): the project own definitions of the reference/slogan/history
% display environments, copied mechanically from the shared Stacks preamble lines 8-12
% (\usepackage{verbatim} plus the three \newenvironment lines, which the chapter preamble keeps
% verbatim from the project preamble). Nothing added, removed or redefined: these environments
% carry no mathematical content and the project itself typesets them as comments. Without these
% three lines a counted statement whose own text uses them fails with Environment slogan undefined.
\usepackage{verbatim}
\newenvironment{reference}{\comment}{\endcomment}
\newenvironment{slogan}{\comment}{\endcomment}
\newenvironment{history}{\comment}{\endcomment}
\usepackage{cleveref}
\crefname{originallemma}{original lemma}{original lemmas}
\Crefname{originallemma}{Original lemma}{Original lemmas}
\makeatletter
\newlabel{algebra-lemma-colimit-finite-etale-given-residue-field}{{09E0}{source}{}{sourcepointer-1}{}}
\newlabel{algebra-lemma-colimit-finite-etale-given-residue-field@cref}{{[originallemma][0][]09E0}{[1][0][]source}{}{}{}}
\makeatother
\crefname{originallemma}{original lemma}{original lemmas}
\Crefname{originallemma}{Original lemma}{Original lemmas}
\makeatletter
\newlabel{algebra-lemma-strictly-henselian-functorial-prepare}{{04GT}{source}{}{sourcepointer-2}{}}
\newlabel{algebra-lemma-strictly-henselian-functorial-prepare@cref}{{[originallemma][0][]04GT}{[1][0][]source}{}{}{}}
\makeatother
\crefname{originallemma}{original lemma}{original lemmas}
\Crefname{originallemma}{Original lemma}{Original lemmas}
\makeatletter
\newlabel{lemma-solutions-go-down}{{34517}{source}{}{sourcepointer-3}{}}
\newlabel{lemma-solutions-go-down@cref}{{[originallemma][0][]34517}{[1][0][]source}{}{}{}}
\makeatother
\title{A separable algebraic base change with separably closed residue fields, in which solutions descend, exists for any extension of discrete valuation rings}
\author{The Stacks Project Authors}
\date{stacks-project:more-algebra.tex:local-snapshot-20261012; GNU FDL 1.2 or later}
\begin{document}
\maketitle
\setcounter{subsection}{1}
\noindent\textit{Editorial note.} All delivered mathematical text is the author's own, copied from the pinned Stacks Project chapter file more-algebra.tex; no statement or proof step was rewritten, reordered or removed.\par\medskip
\noindent\textit{Editorial note.} The chapter file pinned-original-source.tex carries only \textbackslash{}input\{preamble\} on its first line, so no preamble text lives inside the chapter. The retained file sources/stacks-more-algebra/retained-preamble.tex is a mechanical, line-for-line copy of the project's single shared preamble as staged in the snapshot; nothing was added, removed or redefined and no package was added.\par\medskip
\noindent\textit{Editorial note.} The source URL is the source-level Stacks Project URL, as required when more than one card is produced from one source; the per-claim tag is carried by claim\_key.\par\medskip
\noindent\textit{Editorial note.} The counted statement span includes the author's three-line commutative diagram, delivered verbatim. The three references inside the counted proof span are recorded as non-load-bearing navigation pointers.\par\medskip
\noindent\textit{Allowed standard prerequisites (editorial).}\begin{itemize}
\item Solutions, weak solutions and separable solutions for an extension of discrete valuation rings. the property named in the statement at lines 34570-34572. Fixes what is being descended.
\item Filtered colimits of finite etale extensions realizing a prescribed separable algebraic residue field extension. invoked at proof lines 34577-34579. Produces A' with the required residue field and the colimit structure on which the descent argument relies.
\item Strict henselization and the fact that a finite etale algebra over a strictly henselian local ring is a product of copies of that ring. used at proof lines 34585-34598. Is what makes B' tensor\_A A'\_1 split into copies of B', which is the key step of the descent.
\item The solutions-go-down lemma. invoked at proof lines 34604-34605. Carries the solution from A -> B' to A -> B.
\end{itemize}\medskip
\ifdefined\hypertarget\hypertarget{sourcepointer-1}{}\fi
\noindent\textit{Source reference.} Original lemma 09E0: focused-seeds/AGP/sources/stacks-snapshot-20261011/tags.txt, source lines 11883--11883. Cross-chapter navigation pointer into the Stacks Project; the standard fact it supplies is applied inside the counted proof span and the argument resting on it stays there.\par\medskip
\ifdefined\hypertarget\hypertarget{sourcepointer-2}{}\fi
\noindent\textit{Source reference.} Original lemma 04GT: focused-seeds/AGP/sources/stacks-snapshot-20261011/tags.txt, source lines 5844--5844. Cross-chapter navigation pointer into the Stacks Project; the standard fact it supplies is applied inside the counted proof span and the argument resting on it stays there.\par\medskip
\ifdefined\hypertarget\hypertarget{sourcepointer-3}{}\fi
\noindent\textit{Source reference.} Original lemma 34517: sources/stacks-more-algebra/pinned-original-source.tex, source lines 34517--34517. Same-chapter navigation pointer; the statement it names is standard for the section and the argument that consumes it stays inside the counted proof span.\par\medskip
\begin{lemma}
\label{lemma-solution-after-strict-henselization}
Let $A \to B$ be an extension of discrete valuation rings. There exists
a commutative diagram
$$
\xymatrix{
B \ar[r] & B' \\
A \ar[r] \ar[u] & A' \ar[u]
}
$$
of extensions of discrete valuation rings such that
\begin{enumerate}
\item the extensions $K'/K$ and $L'/L$ of fraction fields
are separable algebraic,
\item the residue fields of $A'$ and $B'$ are separable algebraic
closures of the residue fields of $A$ and $B$, and
\item if a solution, weak solution, or separable solution exists for
$A' \to B'$, then a solution, weak solution, or separable solution exists
for $A \to B$.
\end{enumerate}
\end{lemma}

\begin{proof}
By Algebra, Lemma \ref{algebra-lemma-colimit-finite-etale-given-residue-field}
there exists an extension $A \subset A'$ which is a filtered colimit of finite
\'etale extensions such that the residue field of $A'$ is a separable algebraic
closure of the residue field of $A$. Then $A \subset A'$ is an extension of
discrete valuation rings such that the induced extension $K'/K$ of
fraction fields is separable algebraic.

\medskip\noindent
Let $B \subset B'$ be a strict henselization of $B$. Then $B \subset B'$ is an
extension of discrete valuation rings whose fraction field extension is
separable algebraic. By
Algebra, Lemma \ref{algebra-lemma-strictly-henselian-functorial-prepare}
there exists a commutative diagram as in the statement of the lemma.
Parts (1) and (2) of the lemma are clear.

\medskip\noindent
Let $K'_1/K'$ be a (weak) solution for $A' \to B'$. Since $A'$ is a
colimit, we can find a finite \'etale extension $A \subset A_1'$ and a finite
extension $K_1$ of the fraction field $F$ of $A_1'$ such that
$K'_1 = K' \otimes_F K_1$. As $A \subset A_1'$ is finite \'etale
and $B'$ strictly henselian, it follows that $B' \otimes_A A_1'$ is a finite
product of rings isomorphic to $B'$. Hence
$$
L' \otimes_K K_1 = L' \otimes_K F \otimes_F K_1
$$
is a finite product of rings isomorphic to $L' \otimes_{K'} K'_1$.
Thus we see that $K_1/K$ is a (weak) solution for $A \to B'$.
Hence it is also a (weak) solution for $A \to B$ by
Lemma \ref{lemma-solutions-go-down}.
\end{proof}

\end{document}
